\documentclass[11pt,a4paper]{article}
\usepackage[margin=23mm]{geometry}
\usepackage[T1]{fontenc}
\usepackage{lmodern}
\usepackage{booktabs}
\usepackage{xcolor}
\usepackage{hyperref}
\hypersetup{colorlinks=true,urlcolor=blue,pdftitle={Bike Sharing: Forecasting and Planning Decision Brief},pdfauthor={Diogo Ribeiro}}
\setlength{\parindent}{0pt}
\setlength{\parskip}{7pt}
\renewcommand{\arraystretch}{1.15}
\begin{document}
{\LARGE\bfseries Bike Sharing: Forecasts to Decisions}\par
{\large Decision brief \hfill 5 October 2026}\par
Diogo Ribeiro\quad\textbar\quad Reproducible case study, v1.0 preparation

\section*{Decision in one paragraph}
The study supports a transparent discussion of forecast error and planning
preferences. It does \textbf{not support operational deployment on its own}.
The selected calendar model has useful predictive structure, but its uncertainty
bands under-cover the final quarter, especially at high predicted demand.
A cost-aware target illustrates the trade-off between planning too little and too
much. Its costs are assumptions, and the data do not identify actual savings.

\section*{What was evaluated}
The audited dataset contains aggregate hourly rentals in 2011--2012. Four
interpretable candidates were compared in July--September 2012. Forecasts were
issued weekly using only preceding observations. Poisson calendar was selected
by mean absolute error (MAE); signed residual adjustments supplied the planning
policy. October--December was reserved until all choices were frozen.

The final assessment contains 2,168 observed hours and 40 absent hours. Missing
hours were excluded, not filled with zero. Earlier test observations could enter
later scheduled fits and calibration windows; no test outcome selected a model,
policy, cost ratio or method setting. This is sequential evaluation of a fixed
procedure, not one forecast frozen for a whole quarter.

\section*{Point error and uncertainty}
\begin{center}\begin{tabular}{lrrr}
\toprule
Period & Hours & MAE & RMSE \\
\midrule
Validation & 2,208 & 53.268 & 83.215 \\
Final test & 2,168 & 61.378 & 95.811 \\
\bottomrule
\end{tabular}
\end{center}
MAE is the average absolute discrepancy in rentals per hour. Final-test error
is higher than validation error. Only the previously selected model was tested,
so this is not a new competition among all candidates.

\begin{center}\begin{tabular}{lrrr}
\toprule
Nominal & Validation & Final test & Test width \\
\midrule
80\% & 81.25\% & 75.69\% & 150.273 \\
90\% & 90.44\% & 85.19\% & 207.771 \\
95\% & 94.75\% & 91.42\% & 267.424 \\
\bottomrule
\end{tabular}
\end{center}
Coverage is the fraction of observations within the nominal prediction band;
width is measured in rental counts. The nominal 90\% band covers only 85.19\%
overall and 74.81\% in the high predicted-demand group. Recent forecast errors
are therefore not a sufficiently reliable uncertainty model across these
conditions. No method was retuned after seeing this result. The procedure has
no universal or conditional coverage guarantee.

\newpage
\section*{Planning trade-off}
Choose a nonnegative continuous hourly workload target, measured in
rental-equivalent units. Underprediction incurs an assumed cost ratio $r$;
overprediction costs one normalized unit. Loss is
$r\max(y-a,0)+\max(a-y,0)$, for observed rentals $y$ and target $a$.
These are proxy costs, not euros. Fractional targets are allowed because the
quantity is a workload allowance, not a count of bikes or staff.

The validation-selected policy adjusts Poisson forecasts using the
$r/(1+r)$ quantile of past signed standardized errors. It was chosen from five
policies in every scenario. Only those saved choices were tested.
\begin{center}\begin{tabular}{lrrr}
\toprule
Cost ratio & Validation cost & Test cost & Test target \\
\midrule
0.25 & 25.385 & 33.017 & 185.639 \\
0.5 & 36.897 & 45.849 & 205.390 \\
1 & 52.328 & 62.766 & 229.198 \\
2 & 71.368 & 81.273 & 253.693 \\
4 & 90.937 & 102.194 & 275.921 \\
9 & 115.067 & 127.825 & 301.594 \\
\bottomrule
\end{tabular}
\end{center}
Compare periods cautiously: season and observations differ. Compare policies
within one cost scenario, not losses across different cost ratios. Higher
penalties for underprediction produce larger targets and more excess planning.
At ratio 9, targets fall below final-test rentals in 10.10\% of hours; that is an
observed frequency, not a guaranteed service level.

At equal costs the adjusted policy has final loss 62.766, while the separately
frozen Poisson point forecast has MAE 61.378. Thus the small validation advantage
does not persist there. We retain the selected policy and report the adverse
finding rather than select again on the test quarter.

\section*{What would be needed for an operational decision?}
Obtain defensible marginal costs and a mapping from rental-equivalent workload
to actual resources. Station availability, unmet requests, trip movements and
capacity constraints are needed for rebalancing or lost-demand analysis. A design
that measures the effect of changing operations is needed for savings claims.
Observed rentals alone cannot provide that evidence.

The historical notebook explored the full dataset before this modernization.
The maintained code preserves temporal separation, but cannot establish that no
human ever inspected these outcomes. Data lineage, physical weather units and
timezone qualifications remain documented. One two-year dataset is not evidence
of robustness across locations or future operating regimes.

\section*{Evidence and reproducibility}
\href{https://github.com/DiogoRibeiro7/bike-sharing}{github.com/DiogoRibeiro7/bike-sharing}
contains the source, audited provenance, frozen selections, final reports and
exact commands. Tables are generated from committed JSON evidence, not copied
from a spreadsheet. With Python 3.12 and the locked Poetry environment, run
\texttt{poetry run python scripts/reproduce\_release.py} to verify the frozen
assessment. Replays check reproducibility; they do not create another untouched
test or justify tuning to the reported results.
\end{document}
